\section{Topological interpretations}\label{sec:CST}
In this section we give topological interpretations for $T(\xi)$ and $S(\xi)$.
%%%%%%%%%%%%%%%%%%%%%%%%%%%%%%%%%%%%%%%%%%%%%%%%%%%%%%%%%%%%%%%%%%%%%%%%%%%%%%%%%%%%%%%%%
\subsection{Representations}
Let $X$ be the complement of the open tubular neighborhood of $E\subset S^3$.
The fundamental group $\pi_1(X)$ is presented as
\begin{equation*}
  \langle
    x,y
    \mid
    wx=yw
  \rangle,
\end{equation*}
with $w:=xy^{-1}x^{-1}y$.
Let $\rho$ be a non-Abelian representation of $\pi_1(X)$ to $\SL(2;\C)$ given by
\begin{equation}\label{eq:presentation_F_2}
  \rho(x)
  =
  \begin{pmatrix}e^{\xi/2}&1\\0&e^{-\xi/2}\end{pmatrix},
  \quad
  \rho(y)
  =
  \begin{pmatrix}e^{\xi/2}&0\\-d&e^{-\xi/2}\end{pmatrix},
\end{equation}
where $d$ annihilates the Riley polynomial
\begin{equation}\label{eq:Riley}
  d^2-(2\cosh(\xi)-3)d-2\cosh(\xi)+3.
\end{equation}
The preferred longitude $\lambda$ is presented by $y^{-1}xyx^{-2}yxy^{-1}$ and it is sent to
\begin{equation*}
  \rho(\lambda)
  =
  \begin{pmatrix}
    e^{\eta/2}&\ast \\
    0          &e^{-\eta/2}
  \end{pmatrix},
\end{equation*}
where
\begin{multline*}
  \eta
  :=
  \log
  \left(
    \frac{1}{2}
    \left(\vphantom{\frac{1}{2}}
      e^{2\xi}-e^{\xi}-e^{-\xi}+e^{-2\xi}-2
    \right.
  \right.
  \\
  \left.
    \left.
      +(e^{-\xi}-e^{\xi})\sqrt{(e^{\xi}+e^{-\xi}-3)(e^{\xi}+e^{-\xi}+1)}
    \right)
  \right).
\end{multline*}
%%%%%%%%%%%%%%%%%%%%%%%%%%%%%%%%%%%%%%%%%%%%%%%%%%%%%%%%%%%%%%%%%%%%%%%%%%%%%%%%%%%%%%%%%
\subsection{Adjoint Reidemeister torsion}
For a representation $\rho\colon\pi_1(X)\to\SL(2;\C)$, one can consider the cochain complex $C^{\ast}(X;\mathfrak{sl}(2;\C)_{\rho}):=\Hom_{\Z[\pi_1(X)]}(C_{\ast}(\tilde{X};\Z),\mathfrak{sl}(2;\C))$ twisted by the adjoint action of $\rho$.
Here $\tilde{X}$ is the universal cover of $X$, $\pi_1(X)$ acts on $\tilde{X}$ as the deck transformation, and the Lie algebra $\mathfrak{sl}(2;\C)$ is regarded as a $\Z[\pi_1(X)]$-module by the adjoint action of $\rho$.
The Reidemeister torsion $\Tor_{\mu}(\rho)\in\C$ associated with the meridian $\mu$, twisted by the adjoint action of $\rho$ is defined as the torsion of the cochain complex $C^{\ast}(X;\mathfrak{sl}(2;\C)_{\rho})$.
The following formula is known for the case of the figure-eight knot:
\begin{equation*}
  {\Tor}_{\mu}(\rho)
  =
  \pm
  \frac{2}{\sqrt{(e^{\xi}+e^{-\xi}+1)(e^{\xi}+e^{\xi}-3)}}.
\end{equation*}
See \cite{Porti:MAMCAU1997,Dubois:CANMB2006,Murakami:JTOP2013}.
%%%%%%%%%%%%%%%%%%%%%%%%%%%%%%%%%%%%%%%%%%%%%%%%%%%%%%%%%%%%%%%%%%%%%%%%%%%%%%%%%%%%%%%%%
\subsection{Chern--Simons invariant}
Let $M$ be a three-manifold with boundary a torus $T$, and $\{\mu,\lambda\}$ be generators of $\pi_1(T)$.
For a representation $\rho\colon\pi_1(M)\to\SL(2;\C)$, we can define the Chern--Simons invariant as follows.
\par
Let $A$ be an $\sl(2;\C)$-valued $1$-form $A$ on $M$ that defines the flat connection corresponding to $\rho$.
Assume that $\rho\Bigm|_{T}$ is diagonalizable for simplicity.
Then by a suitable conjugation, one has
\begin{equation*}
  \rho(\mu)
  =
  \begin{pmatrix}
    e^{2\pi\i\alpha}& 0               \\
    0               & e^{-2\pi\i\alpha}
  \end{pmatrix},
  \quad
  \rho(\lambda)
  =
  \begin{pmatrix}
    e^{2\pi\i\beta}& 0               \\
    0              &e^{-2\pi\i\beta}
  \end{pmatrix}.
\end{equation*}
Then up to gauge equivalence, we can assume that $A$ is of the form
\begin{equation*}
  \begin{pmatrix}
    \i\alpha& 0       \\
    0       & -\i\alpha
  \end{pmatrix}
  \,dx
  +
  \begin{pmatrix}
    \i\beta& 0      \\
    0      &-\i\beta
  \end{pmatrix}
  \,dy
\end{equation*}
near $T$, where $dx$ and $dy$ are the $1$-forms corresponding to $\mu$ and $\lambda$ respectively.
Then the Chern--Simons invariant $\cs_{M}(\rho;\alpha,\beta)$ of $\rho$ associated with $(\alpha,\beta)$ is defined by
\begin{equation*}
  \CS_{M}(\rho;\alpha,\beta)
  :=
  \frac{-1}{8}
  \int_{M}\Tr\left(dA\wedge A+\frac{2}{3}A\wedge A\wedge A\right)\in\C/(\pi^2\Z).
\end{equation*}
Note that $\Im\CS_{M}(\rho_0;0,0)$ coincides with the hyperbolic volume if the interior of $M$ possesses a complete hyperbolic structure, where $\rho_0$ is the Levi--Civita connection.
See \cite{Kirk/Klassen:COMMP1993} for details.
\par
In \cite{Murakami/Tran:2020}, we proved the following theorem.
\begin{thm}[{\cite{Murakami/Tran:2020}}]
The Chern--Simons invariant $\CS_{E}(\rho;\xi,\eta)$ of the representation $\rho$ associated with $(\xi,\eta)$ is given by
\begin{equation*}
  \CS_{E}(\rho;\xi,\eta)
  =
  S(\xi)-\frac{\xi\eta}{2}.
\end{equation*}
Here we choose $\mu$ and $\lambda$ to be the meridian and the preferred longitude of $E$ respectively, and we assume that $\rho$ sends
\begin{equation*}
  \mu
  \mapsto
  \begin{pmatrix}
    e^{\xi/2}&\ast \\
    0          &e^{-\xi/2}
  \end{pmatrix},
  \quad
  \lambda
  \mapsto
  \begin{pmatrix}
    e^{\eta/2}&\ast \\
    0         &e^{-\eta/2}
  \end{pmatrix}.
\end{equation*}
up to conjugation.
\end{thm}